% STATUS: DRAFT
\chapter{Gravity as constraints}\label{ch:gravity_constraints}

\begin{physmotiv}
In a gauge theory, physical states satisfy Gauss's law. In general relativity, the analogous statement is that physical states are annihilated by the \emph{Hamiltonian and momentum constraints}, the generators of diffeomorphisms. Two consequences drive the rest of the book. First, local fields at a coordinate point are not observables: observables are \emph{relational} (dressed to physical landmarks, such as an observer's clock). Second, in a spatially compact universe there is no Hamiltonian ``at infinity'': the total Hamiltonian is a constraint, so the generator of time translation is itself pure gauge, and the only way to get dynamics is to measure time against a physical clock. The algebra of invariants then turns out to be a crossed product.
\end{physmotiv}

\section{Constraints and gauge invariance}

In canonical gravity the phase space has coordinates $(h_{ij},\pi^{ij})$ on a spatial slice, together with matter fields, and the lapse and shift are Lagrange multipliers enforcing constraints
\begin{equation}
\mathcal H(x)\approx0,\qquad\mathcal H_i(x)\approx0,
\end{equation}
which generate normal deformations and spatial diffeomorphisms of the slice. In the quantum theory one requires $\hat{\mathcal H}\Psi=0$ (Wheeler--DeWitt) and $\hat{\mathcal H}_i\Psi=0$. A gauge-invariant observable is an operator commuting with all constraints. For a closed universe the integrated constraints $\int N\mathcal H$ are the \emph{only} generators of time evolution, so the dynamics is pure gauge.

\begin{keyidea}
Gauge-invariant (Dirac) observables are those commuting with the constraints. A local field $\phi(x)$ at a coordinate point is not one, because moving the point is a gauge transformation.
\end{keyidea}

\section{Relational observables and dressing}

A gauge-invariant version of $\phi$ can be built by \emph{dressing} it to a physical reference: instead of ``the value of $\phi$ at $x$'' one measures ``the value of $\phi$ at the point whose geometric relation to a physical object (observer, clock, field gradient) has given values''. Examples:
\begin{itemize}
\item $\phi$ evaluated along the worldline of an observer, at \emph{proper time $\tau$ read off by the observer's clock}: this is gauge invariant.
\item $\phi$ at the point where a scalar field $\chi$ has value $c$ (if $\chi$ is monotonic): a ``clock field'' in cosmology.
\item In asymptotically AdS or flat spacetimes, operators dressed to the boundary (Gauss-law-like), whose generator is the boundary Hamiltonian.
\end{itemize}
Each choice of reference is a choice of \emph{quantum reference frame} (\cref{ch:clocks_qrf}); different references give different algebras of invariants, and none is canonical.

\section{Group averaging and the physical algebra}

Suppose the constraint generates a group action $G$ on a kinematical Hilbert space $\HH_{\rm kin}$ with unitary rep $U(g)$ and consider an operator algebra $\A_{\rm kin}$ on which $G$ acts by automorphisms $\alpha_g$. Physical states are the $G$-invariant vectors (in the generalized sense of group averaging: $\Psi_{\rm phys}=\int\dd g\,U(g)\Psi$, \cref{ch:groups}) and physical observables are the invariants
\begin{equation}
\A_{\rm phys}=\A_{\rm kin}^{G}=\{a:\alpha_g(a)=a\ \forall g\}.
\end{equation}
When $G=\R$, $\alpha_t=\Ad\ee^{-\ii t\hat G}$ and the kinematical algebra is $\A_{\rm kin}=\M\otimes\Bnd{L^2(\R)}$, with $\M$ a type III algebra and $L^2(\R)$ the Hilbert space of a clock, the invariants are exactly the crossed product of \cref{ch:crossed,ch:toy_clock}. This is how a trace appears: the group average over $\R$ with measure $\dd t$ combines with the weight $\ee^{x}$ to produce $\tau$.

\begin{whatwrong}
\begin{itemize}
\item In a theory with an asymptotic boundary (AdS, asymptotically flat) the Hamiltonian is a boundary term, a non-trivial physical operator and not a constraint, and the algebra of the exterior is a different, generally type III$_1$ or II$_\infty$ object (\cref{ch:bh_crossed}).
\item The constraints of full quantum gravity are not under control beyond perturbation theory. The statements below hold in the limit $G_N\to0$ with the gauge group of \emph{isometries of a fixed background}; this is the sense of ``semiclassical'' throughout.
\item ``Algebra of a region'' is subtle: a region's algebra of dressed operators depends on how the dressing is anchored (to the boundary, to an observer, to the horizon).
\end{itemize}
\end{whatwrong}

\section{Time and the clock: a minimal quantum-mechanical illustration}

A system with Hamiltonian $H_S$ and a clock with Hamiltonian $H_C$ (spectrum $\R$) subject to the constraint $H_S+H_C\approx0$. Physical states: $(H_S+H_C)\Psi=0$. The clock reading $T$ (conjugate to $H_C$) defines ``time $t$ for the system'' via the \emph{relational} observable $a(T=t)=\ee^{\ii H_S t}a\,\ee^{-\ii H_St}$ (Page--Wootters-type construction, for an ideal clock). The invariant algebra is generated by $\{a(T),\,H_C\}$: the crossed product of $\Bnd{\HH_S}$ by the inner flow $\Ad\ee^{\ii H_St}$, which is just $\Bnd{\HH_S}\otimes\Linf(\R)$ in this type I case. For a type III system the flow is outer and the crossed product is genuinely different from a tensor product.

\section*{Exercises}
\begin{exercise}
Verify that $a(T)=\ee^{\ii H_ST}a\,\ee^{-\ii H_ST}$, with $T$ the clock-reading operator ($[T,H_C]=\ii$) commutes with $H_S+H_C$.
\end{exercise}
\begin{exercise}
For a free particle with constraint $H=p^2/2m+H_C\approx0$ show that invariant operators form $\Bnd{\HH_S}\otimes\Linf(\R_{H_C})$ and identify the physical Hilbert space by group averaging $\int\dd t\,\ee^{-\ii tH}$.
\end{exercise}
\begin{exercise}
Explain why in a spatially compact universe with no boundary a gauge-invariant operator cannot be a local field at a fixed coordinate location, and why dressing to a physical reference gives a way out.
\end{exercise}
